\subsection{高一力学基础方法}
\begin{tabularx}{\columnwidth}{XX}
  速度定义 & $v=\frac{\Delta x}{\Delta t}$（极限形式为瞬时速度） \\
  加速度定义 & $a=\frac{\Delta v}{\Delta t}$ \\
  滑动摩擦力 & $f=\mu N$ \\
  牛顿第二定律 & $\sum F=ma$ \\
  二力合力范围 & $|F_1-F_2|\le F\le F_1+F_2$ \\
\end{tabularx}

\begin{formulanotebox}
\begin{itemize}[leftmargin=1.5em]
  \item $x$-$t$ 图像斜率表示速度；$v$-$t$ 图像斜率表示加速度、面积表示位移。
  \item 静摩擦力一般不直接用 $f=\mu N$，通常由平衡或动力学方程反推。
  \item 力学题首要步骤是“选研究对象+画受力图+定正方向”。
\end{itemize}
\end{formulanotebox}

\begin{keypointbox}
高一阶段常见失分点集中在“概念混淆与图像误读”：平均速度和平均速率、位移和路程、受力图漏力或方向错误。
\end{keypointbox}

\begin{casetypebox}[title={\strut\textcolor{white}{\bfseries 常见题型：运动图像判读}}]
\begin{itemize}[leftmargin=1.5em]
  \item $x$-$t$ 图像斜率表示速度，斜率正负表示运动方向，斜率绝对值表示速度大小。
  \item $v$-$t$ 图像斜率表示加速度，图线与时间轴围成的有向面积表示位移。
  \item $v$-$t$ 图像穿过时间轴表示速度反向；图像在时间轴上下的面积可相互抵消位移，但路程要取面积绝对值相加。
\end{itemize}
\end{casetypebox}

\begin{casetypebox}[title={\strut\textcolor{white}{\bfseries 常见题型：受力分析与平衡}}]
\begin{itemize}[leftmargin=1.5em]
  \item 受力顺序按“重力、弹力、摩擦力、其他外力”检查，避免漏力或多画不存在的力。
  \item 静摩擦方向由相对运动趋势决定，大小通常由平衡方程反推，不能直接写成 $\mu N$。
  \item 共点力平衡常用正交分解：$\sum F_x=0$，$\sum F_y=0$。
\end{itemize}
\end{casetypebox}
